\documentclass[letterpaper,11pt]{article}

\usepackage{latexsym}
\usepackage[empty]{fullpage}
\usepackage{titlesec}
\usepackage[usenames,dvipsnames]{color}
\usepackage{verbatim}
\usepackage{enumitem}
\usepackage{fancyhdr}
\usepackage[english]{babel}
\usepackage{tabularx}
\usepackage[hidelinks]{hyperref}

\ifdefined\pdfgentounicode
  \input{glyphtounicode}
  \pdfgentounicode=1
\fi

\pagestyle{fancy}
\fancyhf{}
\fancyfoot{}
\renewcommand{\headrulewidth}{0pt}
\renewcommand{\footrulewidth}{0pt}

% Balanced Margin Adjustments
\addtolength{\oddsidemargin}{-0.6in}
\addtolength{\evensidemargin}{-0.6in}
\addtolength{\textwidth}{1.2in}
\addtolength{\topmargin}{-0.85in}
\addtolength{\textheight}{1.75in}

\urlstyle{same}

\raggedbottom
\raggedright
\setlength{\tabcolsep}{0in}

% Balanced Section Spacing
\titleformat{\section}{
  \vspace{-7pt}\scshape\raggedright\large
}{}{0em}{}[\color{black}\titlerule \vspace{-5.5pt}]

\newcommand{\resumeItem}[1]{
  \item\small{
    {#1}
  }
}

\newcommand{\resumeSubheading}[4]{
  \vspace{-2pt}\item
    \begin{tabular*}{\linewidth}[t]{l@{\extracolsep{\fill}}r}
      \textbf{#1} & #2 \\
      \textit{\small#3} & \textit{\small #4} \\
    \end{tabular*}\vspace{-2pt}
}

\newcommand{\resumeSubSubheading}[2]{
    \item
    \begin{tabular*}{\linewidth}{l@{\extracolsep{\fill}}r}
      \textit{\small#1} & \textit{\small #2} \\
    \end{tabular*}\vspace{-2pt}
}

\newcommand{\resumeProjectHeading}[2]{
    \vspace{-2pt}\item
    \begin{tabular*}{\linewidth}{l@{\extracolsep{\fill}}r}
      \small#1 & #2 \\
    \end{tabular*}\vspace{-3pt}
}

\newcommand{\resumeSubItem}[1]{\resumeItem{#1}\vspace{-2pt}}

\renewcommand\labelitemii{$\vcenter{\hbox{\tiny$\bullet$}}$}

\newcommand{\resumeSubHeadingListStart}{\begin{itemize}[leftmargin=8pt, label={}, topsep=4pt, partopsep=0pt, itemsep=5pt, parsep=0pt]}
\newcommand{\resumeSubHeadingListEnd}{\end{itemize}\vspace{-4pt}}
\newcommand{\resumeItemListStart}{\begin{itemize}[leftmargin=14pt, topsep=2pt, partopsep=0pt, itemsep=1.8pt, parsep=0pt]}
\newcommand{\resumeItemListEnd}{\end{itemize}\vspace{-2pt}}

\begin{document}

\begin{center}
    \textbf{\Huge \scshape Meehir Prabhakar} \\ \vspace{1pt}
    \small +91 6304385747 $|$ \href{mailto:meehir.work@gmail.com}{\underline{meehir.work@gmail.com}} $|$ 
    \href{https://linkedin.com/in/meehir-prabhakar}{\underline{linkedin.com/in/meehir-prabhakar}} $|$
    \href{https://github.com/Meehir5907}{\underline{github.com/Meehir5907}} $|$
    \href{https://meehir.vercel.app}{\underline{meehir.vercel.app}}
\end{center}
\vspace{-12pt}

\section{Education}
  \resumeSubHeadingListStart
    \resumeSubheading
      {VIT-AP University}{Amaravati, India}
      {B.Tech in Computer Science and Engineering}{Aug 2023 -- Jun 2027 (Expected)}
  \resumeSubHeadingListEnd

\section{Experience}
  \resumeSubHeadingListStart

    \resumeSubheading
      {TechtoGreen Drone \& Robotics Private Limited}{Amaravati, Andhra Pradesh, India}
      {AI/ML and Computer Vision Intern}{Dec 2024 -- May 2026}
      \resumeItemListStart
        \resumeItem{\textbf{Autonomous Vision Systems:} Architected a multi-model vision pipeline (YOLO, VGGT, OpenCV) for real-time crop pathology detection across 7 commercial crops and worker safety monitoring.}
        \resumeItem{\textbf{Depth \& Vegetation Analysis:} Combined VGGT monocular depth estimation with YOLO instance segmentation and a custom bitwise mask intersection algorithm to filter foliage within a 1.0-meter proximity threshold.}
        \resumeItem{\textbf{SLAM \& Spatial Computing:} Developed binocular visual SLAM and odometry architectures using StereoSGBM disparity mapping, ORB and FAST feature extraction, and Levenberg-Marquardt pose optimization.}
        \resumeItem{\textbf{Algorithmic Optimization:} Designed custom Non-Maximum Suppression (NMS) algorithms, implemented ArcFace embeddings with FAISS and SQLite for asynchronous POI tracking, and trained Isolation Forests for anomaly detection.}
        \resumeItem{\textbf{Hardware Integration \& Backend:} Synchronized parallel inference with MAVLink telemetry and YDLIDAR G2 point clouds via C++, and built a FastAPI backend generating bilingual English and Telugu PDF agronomic reports.}
      \resumeItemListEnd

    \resumeSubheading
      {The Machine Learning Club, VIT-AP}{Amaravati, Andhra Pradesh, India}
      {Technical \& Project Team Member}{Oct 2023 -- May 2025}
      \resumeItemListStart
        \resumeItem{Built a 3D survival game in Godot (GDScript) featuring state-persistent UIs, dynamic minimaps, and NavigationMesh-based enemy AI pathfinding with progressive difficulty scaling (\href{https://github.com/MLC-VIT-AP/PAC_3D}{\underline{github.com/MLC-VIT-AP/PAC\_3D}}).}
      \resumeItemListEnd

  \resumeSubHeadingListEnd

\section{Projects}
  \resumeSubHeadingListStart
    \resumeProjectHeading
      {\textbf{Dhwani (AR Accessibility Glasses)} $|$ \emph{ESP32-S3, Python, Faster-Whisper, SpeechBrain}}{Aug 2026 -- Nov 2026}
      \resumeItemListStart
        \resumeItem{Architected wearable AR accessibility glasses on an ESP32-S3-WROOM-1 (PlatformIO) with a vertically mirrored SSD1306 OLED holographic HUD (U8g2) and tinted one-way glass.}
        \resumeItem{Integrated a 4-channel INMP441 MEMS microphone array streaming low-latency audio over 2.4 GHz Wi-Fi UDP with custom packet reconstruction.}
        \resumeItem{Built a Python AI pipeline using Faster-Whisper (base.en) with overlap deduplication, GCC-PHAT directional localization, and SpeechBrain ECAPA-TDNN speaker tracking.}
        \resumeItem{Implemented VADER sentiment prefixes and fuzzy name-call detection to trigger real-time HUD subtitles alongside GPIO haptic motor and RGB LED accessibility alerts.}
        \resumeItem{GitHub: \href{https://github.com/Dhwani-AR/ar-accessibility-glasses}{\underline{github.com/Dhwani-AR/ar-accessibility-glasses}}}
      \resumeItemListEnd

    \resumeProjectHeading
      {\textbf{Chakshu (AI SIEM \& MDR Pipeline)} $|$ \emph{Python, Scikit-Learn, SHAP, LIME}}{Jan 2026 -- May 2026}
      \resumeItemListStart
        \resumeItem{Engineered an AI-driven SIEM and MDR framework normalizing heterogeneous telemetry across network flows, web apps, and OS logs into a unified JSON schema.}
        \resumeItem{Deployed Isolation Forests for L3/L4 network anomaly detection and One-Class SVMs with TF-IDF N-grams for L7 web and authentication threat isolation.}
        \resumeItem{Built a central Fusion Engine utilizing a 60-second temporal sliding window for cross-layer cyber threat correlation, integrating SHAP and LIME for explainability.}
        \resumeItem{GitHub: \href{https://github.com/Meehir5907/Chakshu}{\underline{github.com/Meehir5907/Chakshu}}}
      \resumeItemListEnd
  \resumeSubHeadingListEnd

\section{Technical Skills}
 \begin{itemize}[leftmargin=8pt, label={}, topsep=3pt, partopsep=0pt, itemsep=1pt, parsep=0pt]
    \small{\item{
     \textbf{Languages}{: Python, C, C++, SQL, Go, Bash, Kotlin, GDScript} \\
     \textbf{AI, ML \& Vision}{: PyTorch, TensorFlow, Scikit-Learn, OpenCV, YOLO, MediaPipe, FAISS, Transformers, LangChain, SHAP, LIME} \\
     \textbf{Systems \& Tools}{: Linux (KVM, Kernel Tuning), Docker, Kubernetes, Git, GitHub, FastAPI, SQLite, QML, Volatility 3, Wireshark, Suricata} \\
     \textbf{Hardware \& Edge}{: ESP32-S3, Arduino, LiDAR (YDLIDAR), MAVLink Telemetry, Raspberry Pi} \\
     \textbf{Core Domains}{: Computer Vision, Visual SLAM, Cybersecurity (SIEM, MDR), Digital \& Memory Forensics, Unsupervised Learning}
    }}
 \end{itemize}
 \vspace{-4pt}

\section{Awards \& Achievements}
 \begin{itemize}[leftmargin=8pt, label={}, topsep=3pt, partopsep=0pt, itemsep=1pt, parsep=0pt]
    \small{\item{
     \textbf{1st Place, Innovative Model (AI \& DS):} Siddhartha Academy of Higher Education Innovation Day (Oct 2025). \\
     \textbf{AI Hackathon 2025:} AIMERS Letter of Appreciation (Feb 2025). \\
     \textbf{Ideathon 2K25:} Innovation Showcase Participant at Chalapathi Institute of Engineering \& Technology (Mar 2025).
    }}
 \end{itemize}
 \vspace{-4pt}

\end{document}
